% Source: https://artofproblemsolving.com/community/c3h2150563_simple_congruence
% Stable identity: aops-topic-2150563-q01
\textbf{(a)} By definition, $by\equiv1\pmod a$ means that $a$ divides $by-1$. From $ax+by=1$ we get
\[by-1=-ax=a\cdot(-x),\]
which is a multiple of $a$. So $by\equiv1\pmod a$. The argument never used the signs: $a$ divides a number exactly when $-a$ does, so congruence modulo $a$ is the same as congruence modulo $|a|$, and $b$ (or $x$, $y$) may be negative.

\textbf{(b)} Any common divisor of $a$ and $b$ divides $ax+by=1$, so $\gcd(a,b)=1$. Part (a) says $b\cdot y\equiv1\pmod a$, which is exactly the statement that $y$ is an inverse of $b$ modulo $a$.

\textbf{(c)} Here $a=7$, $b=5$, $y=3$, so $5\cdot3=15\equiv1\pmod7$ and the inverse of $5$ modulo $7$ is $\boxed{3}$. (By symmetry, $7\cdot(-2)\equiv1\pmod5$, so $-2\equiv3$ is also the inverse of $7$ modulo $5$.)
